\chapter{Introducción}
\label{cap:introduccion}
\justifying

La transformación digital ha modificado profundamente la manera en que las instituciones administran sus procesos internos, gestionan información sensible y ofrecen servicios a sus usuarios. En este contexto, las organizaciones educativas enfrentan nuevos retos relacionados con la atención de quejas, la protección de los derechos de su comunidad y la necesidad de contar con mecanismos tecnológicos capaces de garantizar seguridad, trazabilidad y eficiencia operativa. Particularmente, dentro del Instituto Politécnico Nacional (IPN), la Defensoría de los Derechos Politécnicos (DDP) desempeña un papel fundamental como órgano encargado de promover, proteger y defender los derechos de estudiantes, docentes, personal administrativo y demás integrantes de la comunidad politécnica.

Actualmente, gran parte de los procedimientos relacionados con la recepción, análisis y seguimiento de quejas continúan dependiendo de procesos manuales o parcialmente digitalizados, lo que genera diversos problemas operativos. Entre las principales limitaciones identificadas se encuentran la saturación administrativa, la dependencia del criterio humano en la etapa de primer contacto, la dificultad para dar seguimiento oportuno a los expedientes y la ausencia de mecanismos inteligentes que permitan automatizar tareas de clasificación y canalización de casos. Estas problemáticas provocan retrasos en la atención, inconsistencias en la interpretación normativa y una limitada capacidad institucional para responder de manera eficiente al incremento constante de solicitudes.

A partir de este escenario surge la necesidad de desarrollar una plataforma web especializada para la gestión digital de quejas y denuncias dentro de la DDP-IPN. La propuesta presentada en este trabajo tiene como objetivo diseñar e implementar una solución tecnológica moderna que permita optimizar el flujo operativo institucional mediante la integración de herramientas de procesamiento de lenguaje natural, aprendizaje automático, arquitectura de microservicios y mecanismos avanzados de seguridad informática. La plataforma busca no únicamente digitalizar formularios o expedientes, sino construir un ecosistema integral capaz de asistir al personal jurídico y administrativo en la toma de decisiones relacionadas con la admisión y clasificación de quejas.

El proyecto se fundamenta en un análisis previo del estado del arte de plataformas gubernamentales, académicas e internacionales orientadas a la gestión de denuncias digitales. Dicho análisis permitió identificar que las soluciones existentes presentan importantes limitaciones relacionadas con escalabilidad, automatización inteligente, protección de identidad, interoperabilidad institucional y especialización universitaria. Mientras algunas plataformas priorizan el cumplimiento normativo y la formalización administrativa, otras se enfocan en anonimización o gestión corporativa, dejando vacíos importantes en contextos educativos donde convergen necesidades jurídicas, administrativas y de protección de derechos humanos. Como resultado, se identificó una oportunidad tecnológica para desarrollar una solución adaptada específicamente a las necesidades operativas de la DDP.

Uno de los principales aportes tecnológicos de este proyecto consiste en la incorporación de un clasificador inteligente de quejas basado en técnicas de Procesamiento de Lenguaje Natural (PLN) y aprendizaje supervisado. Este componente tiene como finalidad asistir al personal de primer contacto en la determinación de competencia o incompetencia de una queja, tomando como referencia el marco normativo establecido en el Acuerdo Número Extraordinario 622 de la Gaceta Politécnica. Actualmente, esta decisión depende del análisis manual realizado por personal administrativo y jurídico, lo cual genera variabilidad en los criterios de interpretación y un consumo considerable de tiempo operativo.

Para resolver este problema, el sistema transforma la narrativa textual de las quejas en representaciones matemáticas mediante técnicas de vectorización y embeddings semánticos. A través de herramientas como spaCy y Sentence-BERT, el modelo es capaz de identificar relaciones contextuales, patrones lingüísticos y características normativas presentes en los textos ingresados por los usuarios. Adicionalmente, se incorporó una arquitectura híbrida que combina información semántica con variables jurídicas estructuradas derivadas del Artículo 4° del acuerdo normativo de la DDP, permitiendo que el sistema no dependa exclusivamente del contenido textual, sino también de reglas legales explícitas.

El desarrollo del clasificador implicó la construcción de un pipeline completo de aprendizaje automático que incluye captura de datos, preprocesamiento lingüístico, generación de embeddings, ingeniería de características, entrenamiento supervisado, validación  e inferencia automática. Para ello se evaluaron distintos modelos de clasificación, entre ellos Máquinas de Soporte Vectorial (SVM), Regresión Logística, Random Forest y modelos de ensamble por votación suave, utilizando métricas especializadas como precisión, recall, F1-score y AUC-ROC. El objetivo principal consistió en minimizar falsos negativos y asegurar que ninguna queja potencialmente competente fuera descartada erróneamente.

Además del componente de inteligencia artificial, el proyecto contempla el desarrollo de una arquitectura moderna basada en microservicios, orientada a garantizar modularidad, escalabilidad y mantenibilidad del sistema. Dentro de los módulos principales destacan el portal del quejoso, el módulo administrativo, el expediente digital, el sistema de seguridad y control de acceso, así como el clasificador inteligente. Esta arquitectura permite desacoplar funcionalidades críticas y facilita futuras ampliaciones del sistema sin afectar la operación general de la plataforma.

Como parte del trabajo realizado, se desarrolló un proceso de interacción constante con la DDP mediante reuniones, levantamiento de requerimientos y validación de procesos operativos. Esto permitió modelar el flujo real de atención de quejas mediante diagramas BPMN, identificar reglas de negocio críticas y diseñar soluciones alineadas con las necesidades institucionales. También se desarrollaron prototipos funcionales, mockups, modelos de base de datos, implementación backend y frontend del microservicio de atención al quejoso, así como pruebas funcionales y experimentales para validar el comportamiento del sistema.

Los resultados obtenidos muestran avances significativos tanto en el desarrollo de la plataforma como en el desempeño del clasificador inteligente. Se logró implementar un sistema funcional capaz de gestionar expedientes digitales, automatizar parcialmente procesos de admisión y ofrecer interfaces modernas para usuarios y personal administrativo. Asimismo, los experimentos realizados evidencian que la combinación de embeddings semánticos y características jurídicas estructuradas mejora considerablemente la capacidad de clasificación del modelo frente a enfoques tradicionales basados únicamente en frecuencia de palabras.

Este trabajo presenta una propuesta integral orientada a modernizar la gestión de quejas y denuncias dentro de la Defensoría de los Derechos Politécnicos del IPN. La integración de tecnologías de inteligencia artificial, arquitecturas escalables y mecanismos de seguridad avanzada permite establecer las bases para una plataforma institucional capaz de optimizar procesos administrativos, reducir tiempos de respuesta y fortalecer la protección de los derechos de la comunidad politécnica. A lo largo de los capítulos siguientes se desarrollan los fundamentos teóricos, metodológicos y técnicos que sustentan el diseño, implementación y validación de la solución propuesta.